\chapter{对国民经济学的接受}

如果我们在近代思想的发展脉络中考察黑格尔与古典政治学和近代自然法的分歧，那么他似乎首先是在实践哲学领域中进一步反对柏拉图—亚里士多德传统，借由理性的批判力量，作别和消解传统的对象与原则。然而，当政治学与自然法的对比限于卢梭《社会契约论》与柏拉图《理想国》的原则上的对照、限于“个别性”和“普遍性”的环节时，真正说来，这算不上是真正具有原创性的对比。如卢梭本人即已进行了这种对比，其尖锐程度远在黑格尔之上。而且卢梭还确实从他所占有的“正确”原则推出了错误的结论：《社会契约论》以柏拉图的《理想国》和亚里士多德的《政治学》为典范，将国家思作“市民共同体”，从而将其等同于“市民社会”(société civile)，由此，其撰构也就重归于传统的外表之下。

与此相反，黑格尔同古典政治和近代自然法的分歧所产生的一个最重要的成果是，他洞见它们的原则的双重不足，由此得出了斩断两者在这个关键点上的传统联系的结论。霍布斯和康德之间的自然法理论设定“个体”意志与“普遍”意志（或者“市民社会”与“国家”）是同一的。与此不同，黑格尔在法国大革命历史经验的影响下，从这样的洞见出发：国家和市民社会的分离应当作为差别的原则得到承认，并应使之在理论上具有效力。古典政治学与近代自然法之间的对照使他认识到，对国家与市民社会的同一性的坚持只是一种假象，同时也与历史上决定两者的前提相冲突。

尽管黑格尔只在《法哲学原理》（1821年）中才论述了现代的、有别于国家的“市民社会”概念，但他在1805-1806年耶拿讲座中即已言及那种“差别”的必然性。由于黑格尔与古典政治学和近代自然法都有分歧，结果这里就为法哲学同时系统地批判柏拉图的《理想国》和卢梭的《社会契约论》做好了准备。近代国民经济学的社会模式使自然法的契约理论遭到了彻底的质疑，黑格尔对近代国民经济学的分析恰好与前述批判形成对照，并弥补了它的不足。

历史上，个体意志与普遍意志的直接同一出现在古代城邦；\footnote{参见《耶拿实在哲学》第2卷，第249页以下：“这是如此令人羡慕并将一直令人羡慕不已的希腊人的美的、幸运的自由。人民同时融入公民之中，人民同时就是个人、政府。它只与自己交往。同一个意志既是个人又是普遍物。然而，必然会出现一种更高的抽象、一种更加伟大的对立和教化，一种更深的精神。”也见罗森克兰茨：《黑格尔传》，第136页。}卢梭（以及青年黑格尔本人）在理论上、罗伯斯庇尔在现实中\footnote{参见《耶拿实在哲学》第2卷，第245页以下。}所做的恢复这种直接同一的努力不仅证明这与从革命释放出来的社会条件不合，而且也与一般现代国家生活的结构形式和要素不合。黑格尔将那种“新科学”，即18世纪产生的、尽管方式不同但与自然法理论一样以个体性原则为基础的国民经济学，也归在一般现代国家生活的结构形式和要素之下。

在当时德国观念论哲学中，黑格尔对从詹姆斯·斯图亚特到亚当·斯密和（在1821年《法哲学原理》中）大卫·李嘉图这些英国经典作家的最先进的国民经济学理论的接受是绝无仅有的。康德基本上将（作为家庭经济、农业和国家经济的）经济理论排除在实践哲学之外，因为它只包含了与社会概念无关的、按照因果性的自然概念才有可能产生结果的技术—实践规则。\footnote{《康德全集》（科学院版）第5卷，第172页。}尽管康德的《法权论》（1797年）第31节接受了亚当·斯密的货币定义（“货币是一种其转让同时是劳动的手段和尺度的物体”），但却鲜明地以个体主义的方式关联到契约上的“‘我的’和‘你的’的转移(commutatio latio sic dicto)中的法概念”，而不是关联到社会的契约建构本身。费希特对国民经济学的看法也是如此。他是重农学派的信徒，比康德还要落后。谢林对这个主题从未表现出更大的兴趣。尽管当时伟大的哲学似乎很少关注这个领域，然而黑格尔对国民经济学的接受也并非孤立事件。从历史上看，它与人们贬为“通俗哲学”的德国启蒙哲学对进步的西欧社会理论的广泛接受浪潮有关。作为新兴的市民社会的科学，国民经济学是那些[启蒙]哲学家们（伽尔韦[Garve]、施洛瑟[Schlosser]、阿普特[Abbt]、伊泽林[Iselin]）想要优先普及的对象。通过他们的介绍，黑格尔熟悉了斯图亚特、斯密、弗格森和休谟的作品。然而，黑格尔由法国大革命及其对德国的影响所产生的实际政治兴趣超过了对公众私人目的进行启蒙的想法。与其18世纪的哲学前辈不同，青年黑格尔研究国民经济学不再出于启蒙公众的需要（在普及科学的过程中，公众确信其私人目的与公共目的实现了同一）；相反，他的研究恰恰是要削弱这种需要，消除那种通过革命进程保证实现这种同一的可能性。在伯尔尼和法兰克福时期，黑格尔试图通过文章和小册子积极介入革命过程，他对现代国民经济学感兴趣的根源似乎就在这里。他认识到，法国大革命对自然法“最佳宪制”理想充满矛盾的实现，其责任并不在于革命的意识形态标准，而在于无法与自然法理想保持一致的社会条件。黑格尔通过分析这些条件，扩展了其革命经验的视野：在法国旁边，矗立着英格兰，这个工业革命的发源地。对黑格尔法哲学的基本历史地位来说，这一点具有决定性的意义。正如卢卡奇中肯地指出的那样，黑格尔不仅“在德国最深刻、最正确地洞见了法国大革命和拿破仑时期的本质，同时还是唯一认真地探讨了英国工业革命问题的德国思想家”。\footnote{G. 卢卡奇：《青年黑格尔和资本主义社会问题》(Der junge Hegel und die Probleme der kapitalistischen Gesellschaft)，柏林，1954年，第25页。}

然而，对于黑格尔法哲学的发展来说，决定性的并不仅仅是那种背景，而且还有这种事实：他将这种由现实的政治兴趣激发起来的对国民经济学的接受从历史的角度与古典政治学和现代自然法关联起来，正因为如此，他就不得不越出它们的界限。诚然，这种关联与青年黑格尔从“应当的东西”发展到认识“存在的东西”的一般转向大致平行，然而他一开始即为这种接受划定了界限。只有在经历多次思想转变后，他才逐渐超越这些界限。本研究将对这些转变进行描述，直到他成功实现了从简单接受国民经济学到阐明现代社会和国家理论的突破为止。与广泛流传的理解——青年黑格尔似乎即已清楚洞见“市民社会的本质”\footnote{参见卡尔·罗森克兰茨：《黑格尔传》，柏林，1844年，第51页。}或者发现他自己面临“资本主义社会的问题”（G. 卢卡奇）——相反，我们坚持认为，直到耶拿时期末，黑格尔都还没有一种允许他把经济研究与政治—历史研究的材料融贯起来的理论。我们手头掌握的，首先不过是这种理论的一些片段，黑格尔处理国民经济学范畴时之所以出现很多矛盾和晦涩，不仅是因为德国社会的落后\footnote{参见G. 卢卡奇：《青年黑格尔》，第127页以下等处。}，而且还因为下面即将谈到的接受过程中一些结构上的特点。

这一过程可以追溯到黑格尔的伯尔尼时期。在其历史—政治的研究和构思的过程中，黑格尔以其特有的细致方式（通过摘录报刊、议会报告等）熟悉了英国的经济和社会状况。根据罗森克兰茨的记载，他对济贫税的辩论尤其感兴趣。\footnote{参见《黑格尔传》，第85页以下。}在对詹姆斯·斯图亚特爵士的《政治经济学原理研究》（1767年）一书德文版（1769年）的评注中，黑格尔将这里获得的对需要和劳动、劳动分工与财富、慈善和赋税（国家管理等）的实践—社会意义的洞见集中地表达了出来。可惜这一评论现已佚失。\footnote{参见上书，第86页，据罗森克兰茨的陈述，[黑格尔的这一评论]写于1799年2月19日至5月16日之间。此外参见夏穆富(P. Chamley)：《斯图亚特和黑格尔的政治经济学和哲学经济学》(Economie politique et philosophique chez Steuart et Hegel)，巴黎，1963年。}除斯图亚特外，在耶拿时期初，亚当·斯密的经济理论又对他产生了重大影响。他通过伽尔韦的翻译（1794-1796年）了解了斯密的《国富论》（1776年）。\footnote{参见《耶拿实在哲学》第1卷，其中黑格尔引用伽尔韦版（布雷斯劳，1794年）第1卷，即斯密。关于讨论黑格尔与国民经济学的分歧的基本文献为G. 卢卡奇：《青年黑格尔》，第208页以下、369页以下。另见K. 赫内(Höhne)：《黑格尔与英国》(Hegel und England)，载：《康德研究》，第36卷，1931年，第301页以下。}黑格尔此时著述的卓异之处在于，他一边研究经济学，一边重构古代城邦理论。就像他在与自然法发生争执时诉诸亚里士多德的《政治学》一样，现在他也以古典经济学原理为背景，处理英国人提出的国民经济学问题。然而，亚里士多德路向的经济学和政治学结构与黑格尔纳入其最早的体系构思当中的那些现实之间的对比之鲜明，远远超过了自然法批判的情形。

在黑格尔最早的政治体系草案《伦理体系》（1802年）中，他与亚里士多德的一致最突出地表现了出来。第一部分（“依照关系的伦理”或“自然伦理”）的关键段落相应于亚里士多德在《政治学》第一卷讨论的那种经济学或家政学(Okonomik)。与在亚里士多德那里一样，黑格尔对人的自然的、出自需要和冲动的活动与团体（劳动、工具的使用、占有、交换；夫妻、父母与子女、主人与奴隶之间的关系）的分析都汇入与自然有关的相对伦理通过支配和家庭经济所能达到的最高“总体”即“家庭”之中：“支配的差别是表面上的。丈夫是主人和管理者，不是与其他家庭成员对立的财产所有者。作为管理者，他只有自由处置的外表。劳动同样按照每个成员的天性进行分配，而产品则是共有的；正是通过这种分工，人人都有剩余，但并不是作为财产。产品无须交换，它直接地、自在自为地就是家庭共有的。”\footnote{参见黑格尔：《政治学与法哲学文集》，第444页。}对自然不同方面（土地、植物、动物）劳动的讨论，都在追踪亚里士多德的主题\footnote{参见《政治学》第1卷，第8章，1256a19—1256b7。}，对具体经济现象的分析（交换、商业、价格、货币）亦然。\footnote{参见《政治学》第1卷，第9章，1257a5—1257b30。另一方面，同样值得注意的是，对亚里士多德暗示的“自动”工具（《政治学》第1卷，第4章，1253b33—1254a1），即机器，黑格尔指出了它的现代形式；他直接从“机械的死板劳动”推演出机器：这种劳动使得“水、风、蒸汽等的运动”作为劳动力进入抽象化的劳动过程成为可能（黑格尔：《政治学与法哲学文集》，第433-434页）。}对黑格尔来说，这个“相对伦理”的所有活动和团体都是通往“人民”这个自足的绝对伦理路途上的前站，而人民则是他按照城邦的模范构建起来的共同体，在这个共同体当中，人类活动的各种可能性都得到了完全的实现。\footnote{黑格尔：《政治学与法哲学文集》，第460页以下。}

然而，黑格尔对古典家政学范畴的部分地、非常紧密的依赖并不能掩盖这一事实，即这些范畴已经再也不能为工业和社会经济解放的事实提供坚实的基础，因为他已经通过研究国民经济学把握了这种解放的意义。在亚里士多德那里，“经济”活动是与“整体”脱离开来的。也就是说，它们要么束缚于一个为同样“个别的”主人而劳动的奴隶的个别性，要么通过其他的方式被排除在作为政治共同体(Koinonia politike)的城邦以外。对黑格尔来说，它们一开始就具有一种特别的“社会的”特性，因为它们是由社会形成的、不依赖于政治的统治活动：“因为劳动以这种方式成为一种普遍的劳动，由此，由于劳动按照其内容不是指向全部需要，相反，按照它的概念，一种普遍的依赖关系便因为自然需要的满足而设定起来。”这样，与亚里士多德的前提完全矛盾，经济学就在《伦理体系》中作为市民的实践，出现在政府的名称之下，表达了公共—政治地位。与在亚里士多德那里不同，在黑格尔这里，经济学取得了这种地位。\footnote{这里，黑格尔即已言及“需要的体系”，这是他后来在《法哲学原理》中分析现代市民社会的基础。然而，我们可以看到，黑格尔尚未将此“体系”设定到“需要”的多样性之中（《法哲学原理》，第三部分，第189-208节），而是从个别“需要”的体系和总体出发的。也见《政治学与法哲学文集》，第490页。}

如果我们追问：黑格尔为什么必须要把经济—政治活动纳入伦理体系之中，那么很快就可以在他那里找到答案。其理由就是：对现代的、自身成为“体系”的政治经济学的接受。在政治经济学中，从需要和劳动产生出来的众多产品的“诸规定性”和“各种实在性”已经获得独立，创造出了一个个人之间的普遍的社会依赖的领域：“在实践领域中，这些处于它们纯粹内在的无形式和单纯性之中的实在性，或者说这些感觉，出于差异，进行自我重建，并且通过消灭诸直观、扬弃了无差异的自我感而得到了重新恢复；——身体的需要和享受重新在总体中自为地设定起来，在它们无限复杂的关系中服从于一种必然性之下，形成了身体需要、劳动以及劳动积累方面的普遍的相互依赖体系，以及——这个体系作为科学——所谓的政治经济学体系。”黑格尔的思想在这个发展阶段必须解决的基本矛盾是，他一方面试图遵照柏拉图—亚里士多德的典范，将人伦政治的（“伦理的”）行动摆到劳动的制造过程之上，同时又要依据现代政治经济学，承认经济制造活动所具有的铸造社会的功能。由此，黑格尔一方面在伦理总体下“纯粹否定地”对待“实在性体系”，让它“接受伦理总体的主宰”，\footnote{参见《政治学与法哲学文集》，第487页以下。}而另一方面，又由于这个体系在其自身得到如此深入的发展，从而“必然摧毁它与那些关系混为一团、完全没有分别开来的自由伦理”。换言之，“有意识地接受这个体系，承认它的权利，在高贵等级之外承认它是一个独特的等级，作为它的领域”。\footnote{参见《政治学与法哲学文集》，第499页。}

从黑格尔最早体系草案对政治经济学体系的这种承认，产生了进一步的矛盾。这个矛盾与上面提到的那个矛盾密切相关。随着将经济—社会领域限制在一个特定等级，对“绝对伦理”的建构也就要从历史哲学的角度进行修正。在黑格尔看来，这是近代世界与古代世界之间的对立所要求的。\footnote{谢林的耶拿时期的《学术研究方法讲座》（1802年夏季学期讲授，1803年印行）也为此提供了动因：“想要作为自由人去理解法和国家的实定科学的每一个人，他的第一个追求必然是：通过哲学和历史获得对晚近世界以及其中的公共生活必然形式的生动直观。”[《谢林全集》第5卷，第315页]也见同书第314页对近代“市民”国家与古代城邦之间区别所作的评论。黑格尔在自然法论文[《全集》第1卷，第497页以下]中的历史哲学论述即由此出发。}然而，这种《伦理体系》尚且陌生的历史哲学修正本身产生了一系列困难。这些困难既与黑格尔建构时从古典政治哲学那里借来的那些材料有关，也与其所接受的近代政治经济学成分有关。对于历史上与近代政治经济学成分相应的那个等级，黑格尔说，它“自身处于一般的占有以及对于占有的可能正义之中，同时，这个等级还构建起一个连贯的体系，并且通过将占有的关系接纳到形式的统一性之中，每一个个人，因其潜在地能够占有，便直接作为普遍物或者作为市民（在作为布尔乔亚[bourgeois]的意义上）与所有人发生关系”。\footnote{《全集》第1卷，第499页。}然而，黑格尔在历史哲学上将这个似乎被他归在近代市民阶层下面\footnote{其更加具体的表现是，他把近代才产生的贵族市民关系也纳入这种论述之中。参见《全集》第1卷，第495页。}的等级解释为古代城邦伦理衰落的产物。随着“绝对伦理的沦丧”，随着城邦的特殊形态消亡于罗马帝国的普遍性之中，古代的自由“等级”和不自由的“等级”成为平等的了，同时“随着自由消失，奴隶制也必然消失”。\footnote{同上书，第496页。}在吉本《罗马帝国衰亡史》的影响下，黑格尔认为，鲜活的城邦伦理僵化为罗马法无生命的形式，作为“私人”之人的普遍性和平等的原则获得了抽象的效力，就像在基督教中，灵魂在上帝面前获得了抽象的平等一样。从历史上看，在一个关键点上，这种说法是错误的（随着人格概念的形成，罗马法绝没有消除黑格尔所言的等级的“分别”，而是在自由人身份(status libertatis)的名义下巩固了这种分别\footnote{参见佐姆(R. Sohm)：《罗马法制度》(Institutionen des Römischen Rechts)，第4版，莱比锡，1891年，第100页以下、105页以下。}）。从其本身来看，如果这种说法不涉及黑格尔如此戏剧性地提出的“伦理悲剧的上演”（这种悲剧在何时何地上演，对这个问题争论不休，绝非偶然）的话，那也并没有大不了的。黑格尔早期著作的解释者们以在其他地方罕见的方式异口同声地对这种悲剧的“晦涩”或“神秘”提出了指控。\footnote{参见卢卡奇：《青年黑格尔》，第465页；H. 格洛克纳：《黑格尔》第2卷，第330页。}实际上，这种“晦涩”或“神秘”的真正原因在于，黑格尔此时对近代“市民社会”的概念和起源还不甚了了。对罗马法的错误解释（这种错误在《法哲学原理》中才得到纠正\footnote{参见《全集》第7卷，第40页以下。}）是他将市民的形象放到古代世界终结处的外部动因，这种投放允许黑格尔在其称赞为所有时代都是有效的埃斯库罗斯的悲剧中去说明历史上随着近代（“市民”）社会的政治和经济解放而产生的作为公民(citoyen)的市民(Bürger)与作为布尔乔亚(bourgeois)的市民之间的冲突。\footnote{参见《全集》第1卷，第500页以下。}伦理悲剧尚未表达出在这种冲突中设立起来的“市民社会”与“国家”之间的差别\footnote{这与G. 洛尔摩泽尔相反。见氏著：《主体性和物化》(Subjektivität und Verdinglichung. Theologie und Gesellschaft im Denken des jungen Hegel, Gütersloh 1961)，第86页以下。他在这一点上追随卢卡奇的主张，大致说来忽略了黑格尔的发展历史。}；相反，它表达的是“第一等级”与“第二等级”之间的差别。前者依照其指向伦理整体的活动，与其自身处于一种无对立的（“无差异的”）关系之中，后者则沉入生产的规定性之中，它的“无差异”只是相对的，受到“现有对立”的束缚。\footnote{参见《全集》第1卷，第500、505页。黑格尔的前提，即绝对伦理通过“贵族和自由人等级”得到表现，更加清楚地出现在1801—1802年第一个等级体系草案中。也见黑格尔：《政治学与法哲学文集》，第472页以下。拉松证实自然法论文已经完成了绝对伦理与第一等级之间的等同（参见《政治学与法哲学文集》，导言，第XXVII页），却又因此而将《伦理体系》的时间放到后面，并且认为其撰构更加成熟，这是不可理解的。}那种“伦理悲剧的上演”，即绝对物所做的永恒的自身游戏，应当保证这种见解：这种对立（第三等级对第一等级的相互依赖以及它们对劳动产品的依赖）就是伦理理念自身的“命运”和“必然性”；因为绝对伦理为了能够在自身深入、巩固的国民经济学体系的历史条件下得到思考，并且作为“形态”得到保持，就必须要这样建构：它通过“牺牲它自己的一部分”，“有意识地承认”一种独立的存在，它也就“从中得到净化，保存了它自己的生命”。\footnote{《全集》第1卷，第537页。}

伦理悲剧的结局，即古代城邦伦理与政治经济学体系的和解，作为该体系通过“分离”和“抵抗”而保存其自身的“命运”，再一次将我们前面讨论过的所有矛盾都保留了下来。这些矛盾构成了黑格尔法哲学和国家哲学第一阶段的特征。它们产生于黑格尔对近代政治经济学和古典政治理论行动与劳动之间关系的同时接受，并且在他引入“伦理体系”结构当中的第一等级与第三等级的独特“分离”中固定下来。显然，真正说来，黑格尔只是在突破了古典樊篱的近代国民经济学基础上重复了政治学和家政学的古典分离，而我们前面讨论过的他的第一个体系和概念构建的困难也与此相连。

其后，黑格尔的思想克服了这些矛盾。当我们从自然法论文的第一次构划来到1803-1804年和1805-1806年耶拿讲座时，就会发现，此时体系和事实似乎完全割裂。基于柏拉图和亚里士多德的划分已经被抛弃。这一时期草稿的特征不仅仅是“材料的任意堆砌”。\footnote{见罗森克兰茨：《黑格尔传》，第194页。}这种堆砌至少可以由此得到部分说明，即这里见到的不是为了出版而写的草稿，而是用于讲座的笔记。而且更加值得注意的是在黑格尔的概念和体系构造的这一阶段所达到的对国民经济学的接受范围，它完全对立于以前的时期。他这一时期阐述的那些范畴既没有受到特殊对象领域的限制，也不固定在一个“等级”上，而是渗入体系的不同部分。其中两点特别有意思：(1)黑格尔将经济学范畴和最近的自然法范畴联系起来，使得政治经济学和（此前一直等同于罗马私法的）“形式”法之间的关系以一种不同的方式出现；(2)他不再将“与需要和劳动有关的相互依赖体系”局限于一个“等级”，而是认识到其中表达了一种关于社会的概念，这个概念使得继续坚持传统的等级学说成为完全不可能的了。遗憾的是，1803-1804年讲座手稿没有对黑格尔政治哲学的这个至关重要的变化为我们提供更多的说明。手稿恰好在黑格尔比较广泛地探讨了国民经济学问题之后，准备向“财产权”和“人格”概念过渡的地方突然中断了。当然，这可能纯属偶然。然而我们只要考虑一下前面段落的处理，就很容易看出，这里，黑格尔消解了其耶拿著作矛盾的一个方面，然而却以极端的形式重复了另一方面。也就是说，一方面，他同先前一样将“绝对伦理”理解为“人民”——个人的存在同它处于纯然否定的关系中；作为“总体”，个体是“一个单纯可能的、并非自为存在的个体，在其现时存在中只是这样的，即个体时刻做好了赴死的准备，舍弃自己。它不是作为个别的总体，作为家庭并且在占有和享受中；相反，这种关系对它自己来说是一种观念性的关系，并且证明为一种自我牺牲的关系”；\footnote{《耶拿实在哲学》第1卷，第231页。}另一方面，黑格尔将政治经济学体系作为绝对伦理的自然基础融入其中。\footnote{参见上书，第234页。在这种联系中，有意思的是，黑格尔在这里将经济活动和政治活动与精神的同样的“必然性”——必须在一族“人民”中工作——联系起来。当然，有别于德性行为的“伦理工作”，第一个等级的工作乃是“否定性的工作，[精神]对准不同于它自己的他者的现象，或是它的无[机]自然”（第234页）。}并且如在自然法论文当中一样，它是伦理的“无机自然”，但同时，它不再只是归于一个“等级”，而是属于“人民”。此前，等级限制了它对于整体的重要性。因此，尽管手稿关于等级学说的段落付之阙如，但是我们还是能够根据黑格尔的论述得出结论。黑格尔在诸如劳动、占有、价值、货币等经济学范畴上面发现的“普遍性”方面本身就是“人民”的定在环节。在这种存在中，这些范畴“直接成为有别于它们自身在其概念中所是的东西”。于是，满足个人需要的劳动——其直接运动形式为个别形式——就转化为一个普遍物：劳动的规则，或是为了所有人的需要的劳动。\footnote{关于劳动的普遍化，参见《耶拿实在哲学》第1卷，第236页：“对劳动本身来说，现在要求同样出现了：[它]要得到承认，具有普遍性形式，正是一种普遍的方式、所有劳动的规则，构成了某种自为存在的东西，表现为一个外在之物、无机自然，并且人们必须要学会它。对劳动来说，这个普遍性是真正的本质。”关于需要的普遍性，参见同书，第237页以下。}在黑格尔看来，这种过程的经典例子就是亚当·斯密所言的劳动分工，制造业中劳动过程的分化。这种分化，作为“劳动的个别化”，无限地增加了生产商品的数量，由此也就无限增加了“个别劳动与完全无限的需要量”之间的相互依赖。现在黑格尔认识到，政治经济学体系在一族“人民”内部构成了一个“团结和相互依赖的体系”，它的“治理”需要不同于等级“区分”和“划界”的其他手段。\footnote{同上书，第239页。另外，也见《论自然法的科学探讨方式》，《全集》第1卷，第498页，这里经济体系的扩张与前述一个等级的构成联在一起。}

1805-1806年讲座对这个体系进行了阐述，它不仅否定了等级区分，而且否定了从自然状态向社会状态过渡的自然法模型。这里，社会劳动的过程取代了自然法模型。在社会劳动过程中，个人将其自然和意识外化为物，以在外化中与其自身进行中介。承认的运动和被承认的存在(Anerkanntsein)，作为普遍的社会中介形式，进入经济过程的运动之中。这份讲稿的伟大在于：一直到费希特为止，现代自然法的社会概念都是抽象的，这份讲稿则通过将现代自然法的社会概念回放到政治经济学的基础上（或者用一个公式来表达：将亚当·斯密和卢梭结合起来），将它从它的抽象中解放了出来。

对于“社会”的创构来说，不仅个人的意志和权利的关系，而且他们在需要和劳动中建立起来的同物的关系，都是构成性的：“被承认的存在是直接的现实性，并且在其要素中，首先就是作为一般自为存在的个人；他享受着、劳动着。唯有这里欲望才有了出场的权利，因为它是现实的了。也就是说，它自身具有了普遍的、精神的存在。所有人为了所有人的劳动，以及享受——所有人的享受。”\footnote{《耶拿实在哲学》，第213页。}对于社会劳动过程来说，物态化(Dingverfassung)或者“意识的自身物化”(Sich-zum-Dinge-machen)\footnote{同上书，第214页：“自我的无机自然，作为存在着的东西，同样与作为抽象的自为存在的自我相对立；自我否定地对待它，并[在]二者的统一[中]扬弃它，然而方式却是这样的，即自我像塑造对象一样塑造其自身，直观它自己的形式，以同样的方式消耗它自己……而这种加工本身就是多方面的；它是意识的自身物化。”}就像个人的劳动一样根本。但在社会中，劳动和需要总是普遍的——“抽象劳动”和“抽象需要”。这种抽象在于“自为存在”、个人在社会中实现的个别化。自为存在和社会两者都是“在其运动中的抽象物”，两者都建立在个人的自为存在着的运动以及个人的劳动之上：“他的劳动的内容超过了他的需要；他为许多人的需要而劳动，每一个人都是如此。因此，每个人都要满足许多人的需要，而他许多特殊需要的满足则来自许多人的劳动。”\footnote{《耶拿实在哲学》第1卷，第214页以下。}劳动和享受过程中的物化同时意味着，物本身进入社会的中介化过程，同时在这些物的上面，被承认的存在的环节、“普遍意志”开始出现。黑格尔在物的经济学存在方式（价值、货币、交换）和法律范畴（如契约）上证明了这一点。在“各种抽象加工”之间必然产生一种“运动”，由此那些为了抽象需要而加工的物便重新成为“具体的[物]”。同时在劳动的抽象中存在着这种运动的可能性，即物与物之间是可以相互比较的。社会劳动是“普遍性要素”，物在其中成为“普遍”和“等值”的了。被加工之物的普遍性就是“它们的等值或价值。在价值上，它们是一样的。这种价值本身作为物，就是货币。回到具体，回到占有，就是交换。在交换中，这个抽象的物表现为其之所是，即从物性(Dingheit)回到自我的变化，并且是以这样的方式：它的物性在于，成为他人的占有……这种物的等值，作为其内在物，就是它的价值，价值完全出自我的同意和他人的意见——我的积极主张，同样还有他的主张，我的意志和他的意志的统一，我的[意志]作为现实的意志、具体存在着的意志而有效；被承认的存在就是定在”。\footnote{同上书，第215页以下。黑格尔边注：“在他的抽象劳动中，他直观到他自身的普遍性、他的形式的普遍性，或者他是为了他人的存在。因此，他想要同他人共享这种设定，或者这些他人自身应当在其中得到直观：第二种运动，它发展并包含了第一种运动的诸环节。自我[是]在与另一个自我的关系中的行动（并且是得到他的承认了的），而另一个自我也关系到我的占有，然而他只有征得我的同意才能拥有我的占有，正如我只有征得他的同意才能[指向]他的占有一样。——两人作为得到承认者的平等。——价值、物的意义。物具有联系到他人的意义。”（第215页以下）}劳动和交换就是将“社会”建构为意志关系和权利关系的中介形式；正如劳动构成了个人与其自身以及作为“物”的自然之间的中介，交换便构成了个人劳动与所有人劳动之间的中介。对黑格尔来说，这些中介形式就是被承认的存在的定在（“事实”），个人在其中直观到其在物的要素中形成的普遍性。在传统自然法中，先占(Okkupation)作为传统自然法财产权的权利资格，在从自然状态向公民状态的过渡中使得社会契约成为必然。在黑格尔这里，与近代国民经济学相一致，先占不再起这种作用。不是先占，而是劳动和交换，构成财产权取得的真正资格，而“社会”的外化形式和中介形式，即被承认的“定在”，则总是构成了财产权取得的前提：“同一个普遍性构成财产权上的中介，作为认知的运动，因此就是直接拥有，它经过了承认的中介，或者说它的定在就是精神的本质。这里，据有的偶然性被扬弃了：通过劳动和交换，我在被承认的存在中拥有了一切。……这里，财产权的渊源和起源是劳动、我的活动本身的这种[普遍性]——直接的自我和被承认的存在和根据。”\footnote{《耶拿实在哲学》第1卷，第217页。参见同页对这一部分的总结。这里首先需要注意的是，黑格尔不仅将交换和契约而且将劳动理解为外化(alienatio)的形式：“我已经在交换中表达意志，将我的物设定为价值。也就是说，内在的运动、内在的活动，正如劳动是沉入存在之中的活动一样；[在两种情况下都是]同一种外化。（在劳动中，我使自己直接成为物，[成为]形式，即存在。）以同样的方式，我外化我的这个定在，使之成为我的异己的存在，并在其中保存自我……因此，这里，我就将我的得到承认的存在直观为定在，并且我的意志就是这种效力。”（第217页）}

由此，黑格尔将劳动理解为现代社会的解放形式，在现代社会，个人被“教化”为法律人格的自由。无论是在古典政治学中，还是在近代自然法中，劳动的这种“社会”性都隐而未显：前者将伦理行动置于制造之上；后者则着眼于意志关系和契约关系。[而在黑格尔这里，]契约的“社会”事实则是这种本源性的承认(Anerkanntsein)的结果，个人通过劳动和交换而获得这种承认。这种确定性，即个人之间的关系和交付通过契约而获得保证，使得建立在劳动和交换之上的社会的偶然关系成为可预测的和理性的了：“我的话必定有效，这不是出于我内心诚实，我的意向、信念等不可更改的道德理由”，而是因为意志只有作为“被承认的[意志]”或者“社会的[意志]”才有其定在：“我不仅自相矛盾，而且与我的意志要得到承认这一点相矛盾。人们不能相信我的话。也就是说，我的意志只是我的，只是单纯的意见。”\footnote{参见《耶拿实在哲学》第1卷，第219页以下。此外，参见H. 马尔库塞：《理性和革命》(Vernunft und Revolution, Neuwied, 1962)，第81页以下。}契约把个人当作自由、平等的个人对待，个体意志不是个别的意见，而是可以预见的普遍意志。诚然，在契约（黑格尔将契约称为“观念性的交换”、主张的交换，而非物的交换）中，“概念”和“定在”彼此分离；由于个体坚持其个体意志，而非普遍意志，这样就可能会违反契约。只有履约才是“定在”或者“存在着的普遍意志”。\footnote{同上书，第218页以下。}

因此在契约中，普遍意志也就“隐藏在特定的物中”\footnote{参见上书，第220页。}，这样，个人便在外化的行为中彼此结合起来。在黑格尔看来，普遍意志必须要走出这种隐藏状态；因为契约很好地规范了整个社会内部个人之间的关系，但是并没有规范社会内部的运动。这里，在对这种运动之规律的认识中，黑格尔完全超出了自然法的视域。他在“被承认的存在”这一部分（这一部分至少形式上坚持了从自然状态到公民状态过渡的理论）之后，在讨论“伦理”（政治学，这里包含了国家宪法和行政理论[宪政]）之前，进一步论述了“掌握权力的法律”，在法治的视角下对前面讨论过的材料进行了重新处理。这里很容易看到，现在，政治经济学在何种程度上更加广泛、系统地进入政治哲学当中。\footnote{1803-1804年讲座绝没有讲得这般清楚。尽管黑格尔在这里讨论“绝对伦理”时重复了语言、工具、家庭、家产等“自然阶次”，但是他并未改变“人民”伦理的体系地位。相应地，作为第一部分之结束的承认过程，并不以一个独特的“得到承认的存在”的领域为结果，而是直接过渡到对伦理的阐述。1805-1806年的“得到承认的存在”构成“现实精神”的第一阶段(甲)，它与(乙)契约、(丙)犯罪与惩罚和(丁)掌握权力的法律一道，构成了体系的第二部分，并且作为独立的领域，插入体系的第一部分（理论自我和实践自我的发展史）和第三部分（宪法，即国家理论和行政理论）之间。对此，参见F. 罗森茨威格：《黑格尔与国家》第1卷，第178页以下。事实上，他并未对（紧随耶拿早期论文之后的）1803-1804年讲座体系做出中肯的判断。}黑格尔在前面的部分将“普遍性要素”描述为经济法律范畴的“社会”环节，此时这些范畴还带有个人行为所具有的偶然性，这里，他明确地直呼其名：现在，正是社会本身在“掌握权力的法律”的标题下找到了它的位置：“普遍物……对个别劳动者来说是一种必然性。他在普遍物上面有其无意识的存在；社会是他的本性，他依赖于社会的根本的、盲目的运动，这种运动在精神上和肉体上维护他，或是毁灭他。”\footnote{《耶拿实在哲学》第1卷，第231页。}政治上得到解放、按经济规律运行的“社会”以双重方式出现：(1)以肯定自为存在的个人的偶然性和权利的形式，由此社会将自身作为整体重新产生出来，并且表明自己优越于传统政治社会的实体性伦理。这种肯定的表达就是黑格尔所言的将人格、“个人的自我包含于自身”的法律。\footnote{《耶拿实在哲学》第1卷，第225页。参见同书第226页以下对法律各环节的描述。}然而，(2)这个社会处于“自然必然性、个人的偶然性”之下，以至于当它独立自为时，就会产生许多对立，使它和个人的定在面临碎裂的危险。个人的劳动，维持其存在的可能性，“处于完全交织在一起的整体的偶然性之下。因此大量的人口注定要在工厂、制造业和矿山中从事完全机械的、不健康的、不安全的低技术水平的劳动。由于外国发明等造成的[新的生产]方式或者改良，养活了一个大的阶级的工业部门一夜之间便垮掉了，所有这些人随之陷入他们自己无能为力的贫困之中。巨贫与巨富之间的对立出现……万般无奈的贫困”。\footnote{同上书，第232页。}

尽管这种对立摧毁了“人民的绝对纽带、伦理”，但《伦理体系》仍然“直接通过等级自身的构成”来保证整体的维系。\footnote{参见黑格尔：《政治学与法哲学文集》，第492页。尽管对立产生于“需要的体系”内部，但因为黑格尔此时在其等级学说的视角下思考“需要的体系”，因此对立便出现在“营利等级”内部。盖营利等级“又重新分为许多特殊的营利等级，并将这些等级区分为不同财富和享受的等级”（第491页）。关于政府强加的外部限制，参见同书，第493页以下。与等级的建构相比，这一部分的论述相对简略，基本上只涉及税务制度。}1803-1804年讲座不再将这个对立简化为一个特殊等级，而是将其理解为那个设定在“人民”自身中的“巨大的团结和相互依赖的体系”的一个环节。但由于讲座在讲到这个体系时中断了，\footnote{参见《耶拿实在哲学》第1卷，第239页以下。}因此对此问题我们也就不得而知了。1805-1806年讲座则有所不同，它在社会的现代名义之下，将该体系作为“国家”与人民的“宪法”外在地分开。它不仅花了大量的篇幅分析社会的对立，\footnote{参见《耶拿实在哲学》第2卷，第231-233页。}而且还比之前更纯粹、更清楚地论述了这些社会对立的“必然性”的辩证法。社会整体的根本必然性渗入社会个体的意识之中，同样也关联到偶然性的基本形式。通过偶然的运动，社会整体维持其自身和个人。这一法则在肯定个人偶然性的同时，又必然再度否定之。也就是说，社会的盲目运动自己摆脱偶然性，使它自己的定在归于毁灭。对立的社会本身无法摆脱这个魔咒；为此就需要一个救急神(deus ex machina)，在法则中将必然和偶然结合起来。在黑格尔看来，这个平衡社会对立并使整个社会运动与其自身和个人的运动实现和解的法则，就是国家：“然而，这种构成个人存在的完全偶然性的必然性，同样是维持个人存在的实体。国家权力进来，并且必然关心每个领域的维持，或者进行干预，在海外寻找出路、销售渠道等……职业自由[是必需的]，干预必须尽量不要表现出来，因其属于任意的领域。”\footnote{同上书，第233页。}因此，在黑格尔从对“社会”的经济运动的论述过渡到国家（三、宪法）之前，国家便出现在“掌握权力的法律”的概念之下。尽管黑格尔明确提到了社会的名称，个人的“法律”也已经区分为家庭（“个人直接存在的持存”\footnote{参见上书，第227页：“A. 法律[作为]其直接定在的持存。个人在其中直接作为自然的整体，个人作为家庭。他作为这个自然整体，不是作为个人(Person)；他必须首先成为个人。”另见同书第228页以下。}）的法律和家庭解体后的人格与财产领域的法律，然而在其思想的这一阶段，他仅将处于家庭和国家之间的个人存在叫作“国家”。\footnote{《耶拿实在哲学》第2卷，第231页：“B. 这种个人直接定在的法律，作为法律，构成他的意志，或者它在偶然存在的消失中维持此意志本身。[通过]父母离世，它成为肯定的；它作为父母之前所是的定在出现——国家。”另见罗森茨威格：《黑格尔与国家》第1卷，第180页以下。}

由此，这里概念和得到表现的实在性仍然是分开的。但同时概念（“国家”）越来越具有了讲座大量涉及的材料的历史外貌。在最后的段落，国家是“作为财富的国家”，而在“宪法”部分，国家成为了作为“普遍力量”和“绝对权力”的国家。\footnote{参见上书，第242、244页以下。}黑格尔似乎在修改这一部分时利用了1799至1802年间所写的关于德国宪法的手稿。这份手稿的第一部分提出的“国家的概念”在许多方面同讲座的国家概念是一样的\footnote{见黑格尔：《政治学与法哲学文集》，第17页以下；《耶拿实在哲学》第2卷，第246页以下。这种一致首先来自主权的特征以及自保的规则方面，为此黑格尔在这两部作品中都引用了马基雅维利的《君主论》。——海姆在《黑格尔及其时代》第167页正确地主张，《德国宪法》中的国家概念差不多完全跟早期哲学草案中“构建的”国家概念矛盾。然而，海姆没有提到这一事实，即这种看法并不适用于1805-1806年讲座。}；无疑，黑格尔的“国家”同现代政治概念（从马基雅维利到霍布斯阐发的自保和权力行使的规则）之间的联系，与前面观察到的进一步将国民经济学纳入政治哲学之中，是联系在一起的。由此而在政治学和自然法的传统概念和体系架构中构建起来的现代社会概念，一方面使人再也不能死抱住传统等级学说不放，另一方面也使得国家概念再也不能片面定位于古代的“实体性伦理”。这样，黑格尔便在本次讲座中最终与其青年时代的城邦理想诀别。他看出了这个理想与社会契约理论之间的隐秘联系：“有人如此设想普遍意志的构成：所有公民都参加集会、商议、投票，于是多数便形成了普遍意志。”\footnote{《耶拿实在哲学》第2卷，第245页。}黑格尔指责自然法理论，因其假定，普遍意志唯有出自个体意志方能形成，而另一方面又预设，甚至在人们集会成为国家之前，普遍意志即已存在于他们心中。在黑格尔看来，这种契约模式与“国家”（这里理解的国家为“近代”国家，既包括绝对主义国家，也包括革命以后的国家）的历史和社会基础不合。讲座非常清楚地指出了这一点。普遍意志不是国家的原因，而是国家的结果。国家的起源来自外部力量：“国家是他们的自在的本质(Ansich)，也就是说，是对他们进行强制的外部力量。因此，所有国家都是由伟大人物的非凡力量缔造的……”\footnote{同上书，第246页。}由此，黑格尔转向其时代的国家，对他来说，就是被法国大革命动摇了的、其基础得到了重新改造的近代欧洲国家。革命实现了社会契约，社会契约认为，普遍意志由所有人的意志构成，它是“真正的、自由的国家的原则”。\footnote{参见上书，第245页。}然而革命并没有实现城邦理想，而是实现了其在近代国家的经济社会条件下已经有所不同的摹本。契约理论的出发点是普遍意志与个体意志之间的直接同一。从历史的现实性来看，这种直接统一已经分裂为“自身即为个体性、[作为]政府的普遍物的端项，[政府]不是一个国家的抽象物，[而是]以普遍物自身为目的的个体性，以及以个人为目的的另一端。而这两种个体性是一样的。同一个人既关心自己和家庭；他劳动、订立契约等，同时也为普遍物劳动，以普遍物为目的。按前一方面，他叫市民(bourgeois)；按后一方面，他叫公民(citoyen)。所有人都要服从作为多数意志的普遍意志，[它]通过众多个人的特定表达和投票形成”。\footnote{《耶拿实在哲学》第2卷，第249页。黑格尔边注：“小市民(Spießbürger)和帝国国民(Reichbürger)，两者形式上都是市民。”}

个体性和普遍物的统一在历史上得以“定”在的“双重方式”已经暗示了1821年《法哲学原理》的结构：市民社会理论（人作为市民[bourgeois]）与国家理论（人作为公民[citoyen]）的分离。然而，此时（1805-1806年）黑格尔自己还对这个由法国革命提出的人[homme]作为公民与人作为市民的区分问题并不是很清楚。上述引文之后，黑格尔仍然直接过渡到对古代民主的描述，从而我们也就无从知道，法国的、为民主的革命实现奠定基础的这个对立与“希腊人美的、幸福的自由”究竟有何干系。然而，讲座对一点，即现代国家优越于古代“民主”的原则，是给出了明确答复的。这里显然，正是近代国家的原则自身产生了市民和公民的对立，并在思想和现实中将这种对立推向顶点，由此而能够承受它：“因此，更高的分裂就是，每个人都完全回到自身，知道他自己本身就是本质，以至于坚持自己与定在的普遍物分离，而绝对存在。”在黑格尔看来，近代国家的原则在于，个人让普遍物获得自由，普遍物也同样让个人获得自由。通过彼此都让对方获得自由，由此两个领域都获得了独立。普遍意志不仅由个体意志构成，而且还摆脱了个体意志——“摆脱了所有人的知识”。\footnote{同上书，第249页以下。}黑格尔由此推出了世袭君主制是与近代世界相符的国家形式的结论。从法哲学和国家哲学的角度来看，这个结论并不像如下事实那样具有决定性意义：他用这种原则摧毁了从契约推出国家概念的基础，由此也就摧毁了政治（“市民”）社会概念的基础。现在，随着同“社会”的分离，国家才首次获得了一个独特的、与其现代概念相符的地位，并成为后来法哲学体系的出发点。国家在两方面获得了自由：相对于个体意志的自由，个人落入“市民”独立性的要求之中；相对于君主属性的自由（君主只是一个“组织起来的”整体的“空洞的节点”，整体则生活于制度之中，并且知道承受它在其内部形成的那些对立[“极端”]）。\footnote{参见《耶拿实在哲学》第2卷，第251页：“现在，在那些宪法依旧的国家，统治和生活已经发生改变——这是随着时间逐渐改变的……每一个领域：城市、行会在管理其特殊事务时[都是实际的]。但人民构成政府时，就是糟糕的——像非理性一样糟糕。然而整体就是中心，就是自由的精神，它摆脱了那些完全固定的极端，也摆脱了君主的个性；[他]是空洞的节点。”另见同书，第252页。类似思想已经出现在《论德国宪法》（1801-1802年）中，载黑格尔：《政治学与法哲学文集》，第27页。}